% ============================================================
%  Chapter 5 — Tiny Cameras, Big Pictures
%  Inside a Smartphone: The Tiny World in Your Hand
% ============================================================

\chapter{Tiny Cameras, Big Pictures}

\chapteropening{%
  One small hole in the back of your phone can capture a moment forever.
  How does such a tiny camera take such sharp, colourful photos?
}
\chapterbadge{Mission: Capture the World with Light!}

% --- Opening ---
\section*{Meet Lens}

\character{Lens}{Hello! I am Lens — and light is my favourite thing.
  I bend it, focus it, and send it exactly where it needs to go.
  Without me, every photo would be a blurry mess.
  Together with the sensor behind me, we turn a moment into a memory.}

\sectionrule

% --- How cameras work ---
\section*{Lenses and Sensors}

Every smartphone camera has two main parts working together.

\subsection*{The Lens}
A \term{lens} is a curved piece of glass or plastic.
When light enters the camera, the lens bends it so that it focuses
into a sharp, clear image on the sensor behind it.
The lens helps light from each tiny part of the scene land in the right place,
so details do not smear together.

Without a lens, the light would scatter and produce a blurry image.

\subsection*{The Camera Sensor}
Behind the lens sits the \term{camera sensor}.
It is covered in millions of tiny light-detecting spots called \term{pixels}.
When light hits the sensor, each pixel records how bright the light is
and what colour it is.
All those tiny measurements combine to form a complete photograph.
Then the processor quickly combines all that information into the photo you see.

\begin{picturethis}
  Imagine holding a piece of graph paper up to a colourful painting.
  Each square on the grid captures one tiny patch of colour.
  The more squares there are, the more detail you get.
  A camera sensor works the same way — each pixel records one tiny
  patch of the scene.
\end{picturethis}

\begin{funfact}
  A modern smartphone camera sensor can have \textbf{50 megapixels} —
  that is 50 \emph{million} tiny light detectors, all in a chip smaller
  than your thumbnail.
\end{funfact}

\sectionrule

% --- Selfies ---
\section*{How Selfies Work}

Most smartphones have \textbf{two cameras}:
\begin{itemize}
  \item The \term{rear camera} (on the back) — usually higher quality, used for general photos.
  \item The \term{front camera} (on the screen side) — used for selfies and video calls.
\end{itemize}

When you tap the camera-flip button, the phone simply switches which sensor
is active. The front camera works in exactly the same way as the rear one.

\sectionrule

% --- Night mode and zoom ---
\section*{Night Mode and Zoom}

\subsection*{Night Mode}
In low light, a camera sensor does not receive much light per pixel,
making photos look dark or grainy.

\term{Night mode} solves this by taking many photos very quickly
and combining them into one image. The processor averages out the
light information across all the shots, revealing details invisible in any single photo.

\subsection*{Zoom}
There are two types of zoom:
\begin{itemize}
  \item \term{Optical zoom}: the lens physically moves to magnify the image.
        This keeps the photo sharp.
  \item \term{Digital zoom}: the processor crops and enlarges the image.
        This can make the photo look less sharp.
\end{itemize}

Many smartphones have two or three cameras with different focal lengths,
so they can switch lenses to get true optical zoom.

\begin{diagram}[How light travels through a lens to the sensor.]
  \begin{tikzpicture}[>=Stealth]
    % Light rays coming in from left
    \draw[WarmOrange!70,thick,->] (-3.5, 0.8) -- (-1.2, 0.8);
    \draw[WarmOrange!70,thick,->] (-3.5, 0.0) -- (-1.2, 0.0);
    \draw[WarmOrange!70,thick,->] (-3.5,-0.8) -- (-1.2,-0.8);
    % Lens
    \fill[SkyBlue!40] (-1.2,-1.0) .. controls (-0.5,-0.6) and (-0.5,0.6) .. (-1.2,1.0)
                   -- (-0.9,1.0) .. controls (-0.2,0.6) and (-0.2,-0.6) .. (-0.9,-1.0) -- cycle;
    \draw[SmartBlue,thick] (-1.2,-1.0) .. controls (-0.5,-0.6) and (-0.5,0.6) .. (-1.2,1.0);
    \draw[SmartBlue,thick] (-0.9,-1.0) .. controls (-0.2,-0.6) and (-0.2,0.6) .. (-0.9,1.0);
    \node[below,font=\small] at (-1.05,-1.2) {Lens};
    % Rays converge to sensor
    \draw[WarmOrange!70,thick,->] (-0.9, 0.8) -- (1.4, 0.0);
    \draw[WarmOrange!70,thick,->] (-0.9, 0.0) -- (1.4, 0.0);
    \draw[WarmOrange!70,thick,->] (-0.9,-0.8) -- (1.4, 0.0);
    % Sensor
    \fill[TechPurple!60] (1.4,-1.0) rectangle (1.8,1.0);
    \node[right,font=\small] at (1.85,0) {Sensor};
  \end{tikzpicture}
\end{diagram}

\sectionrule

\begin{newwords}
  \begin{description}[labelwidth=3.5cm, leftmargin=3.7cm]
    \item[\term{Lens}]           A curved piece of glass that focuses light.
    \item[\term{Camera sensor}]  The chip that converts light into digital information.
    \item[\term{Pixel}]          One tiny light-detecting spot on a sensor or screen.
    \item[\term{Night mode}]     A feature that combines multiple shots for brighter low-light photos.
    \item[\term{Optical zoom}]   Zoom achieved by moving the lens, keeping images sharp.
    \item[\term{Digital zoom}]   Zoom achieved by cropping and enlarging the image.
    \item[\term{Selfie}]         A photo taken of yourself with the front-facing camera.
  \end{description}
\end{newwords}

\sectionrule

\begin{experiment}
  \textbf{Light and shadow!}

  Take a photo of the same object in three conditions:
  \begin{enumerate}[label=\textbf{\arabic*.}]
    \item In bright daylight or near a window.
    \item In a dimly lit room.
    \item In a dimly lit room with night mode turned on.
  \end{enumerate}
  Compare the three photos. What differences do you notice in sharpness,
  brightness, and colour? Can you explain why?
\end{experiment}

\sectionrule

\begin{bigidea}
  A smartphone camera uses a \textbf{lens} to focus light onto a tiny
  \textbf{sensor} packed with millions of pixels.
  Smart software like night mode helps the processor
  produce clear photos even in difficult conditions.
\end{bigidea}

\character{Lens}{A great shot needs power too. Let’s visit Zap
  and learn about the battery in Chapter 6!}
